\documentclass[a4paper,12pt]{article}

\usepackage[T2A]{fontenc}
\usepackage[utf8]{inputenc}
\usepackage[russian]{babel}
\usepackage{amsmath,amssymb,mathtools}
\usepackage{helvet}
\renewcommand{\familydefault}{\sfdefault}
\usepackage{icomma}
\usepackage[a4paper,margin=2.05cm,headheight=25pt]{geometry}
\usepackage{microtype}
\usepackage{tikz}
\usepackage{tabularx,booktabs}
\usepackage{enumitem}
\usepackage{parskip}
\usepackage{fancyhdr}
\usepackage{setspace}
\usepackage{xcolor}

\definecolor{accent}{HTML}{008080}
\definecolor{darkaccent}{HTML}{004D4D}
\definecolor{textgray}{HTML}{333333}
\definecolor{softgray}{HTML}{6B7280}
\definecolor{verylight}{HTML}{EEF4F4}

\onehalfspacing
\setlength{\parindent}{0pt}
\setlist[itemize]{leftmargin=1.5em,itemsep=0.25em,topsep=0.35em}
\setlist[enumerate]{leftmargin=1.8em,itemsep=0.45em,topsep=0.35em}
\renewcommand{\arraystretch}{1.3}

\newcolumntype{L}[1]{>{\raggedright\arraybackslash}p{#1}}
\newcolumntype{C}[1]{>{\centering\arraybackslash}p{#1}}
\newcolumntype{Y}{>{\raggedright\arraybackslash}X}

\pagestyle{fancy}
\fancyhf{}
\fancyhead[L]{\small\textcolor{softgray}{Классическая вероятность: табличный метод}}
\fancyhead[R]{\small\textcolor{softgray}{Ксения Клыкова}}
\fancyfoot[C]{\small\textcolor{softgray}{\thepage}}
\renewcommand{\headrulewidth}{0.4pt}
\renewcommand{\headrule}{%
  \hbox to\headwidth{%
    \color{accent!55}\leaders\hrule height \headrulewidth\hfill
  }%
}

\newcommand{\sectionline}{%
  \vspace{-0.3em}%
  \textcolor{accent!65}{\rule{\linewidth}{0.6pt}}%
  \vspace{0.4em}%
}
\newcommand{\key}[1]{\textbf{\textcolor{darkaccent}{#1}}}
\newcommand{\badcell}{\textcolor{softgray}{$\times$}}
\newcommand{\goodcell}[1]{\textbf{\textcolor{darkaccent}{#1}}}

\begin{document}
\color{textgray}

\begin{titlepage}
    \thispagestyle{empty}
    \centering
    \vspace*{2.1cm}

    {\Large\textcolor{darkaccent}{Чек-лист}\par}
    \vspace{0.75cm}

    {\Huge\bfseries Классическая вероятность\par}
    \vspace{0.35cm}

    {\LARGE Табличный метод\par}

    \vspace{1.1cm}
    {\large ЕГЭ по математике\par}

    \vfill

    {\large Лёвин Артём Александрович\par}
    \vspace{0.2cm}
    {\large эксклюзивно для Ксении Клыковой\par}
    \vspace{0.7cm}
    {\normalsize \today\par}

    \vfill

    \begin{tikzpicture}[x=1cm,y=1cm]
        \draw[accent!55,line width=1.2pt]
        (0,0)
        .. controls (2.1,0.52) and (4.7,-0.23) .. (6.9,0.10)
        .. controls (9.0,0.40) and (11.1,0.14) .. (13.2,0.34);
    \end{tikzpicture}

    \vspace*{0.7cm}
\end{titlepage}

\section*{1. Главная формула}
\sectionline

Если все элементарные исходы равновозможны, то
\[
P(A)=\frac{m}{n},
\]
где $n$ --- число всех допустимых исходов, а $m$ --- число исходов,
благоприятных событию $A$.

На этом этапе изучения вероятности количество исходов определяем
\key{табличным методом}: строим таблицу, выписываем в её клетках все
допустимые варианты и только после этого считаем клетки.

\key{Два главных вопроса перед построением любой таблицы:}
\begin{enumerate}
    \item Допускаются ли дубли, то есть может ли один и тот же объект
    или результат встретиться дважды?
    \item Важна ли очередность, то есть различаются ли, например,
    $AB$ и $BA$?
\end{enumerate}

Ответы на эти вопросы определяют, какие клетки таблицы допустимы
и нужно ли учитывать зеркальные клетки дважды.

\section*{2. Универсальный алгоритм табличного метода}
\sectionline

\begin{enumerate}
    \item
    \key{Определите смысл строки и столбца.}

    Например: строка --- выбор первого человека, столбец --- выбор второго.

    \item
    \key{Запишите все возможные значения} по строкам и столбцам.

    \item
    \key{Заполните клетки исходами.}

    Одна клетка должна соответствовать одному полностью определённому
    результату опыта.

    \item
    \key{Уберите недопустимые клетки.}

    Если повтор запрещён, клетки вида $AA$, $BB$, $CC$ не учитываются.

    \item
    \key{Проверьте очередность.}

    Если $AB$ и $BA$ описывают один и тот же результат, учитывайте только
    одну из двух зеркальных клеток.

    Если $AB$ и $BA$ являются разными результатами, учитывайте обе клетки.

    \item
    \key{Отметьте благоприятные клетки} и посчитайте
    \[
    n=\text{число всех допустимых клеток},
    \qquad
    m=\text{число благоприятных клеток}.
    \]

    \item
    Найдите вероятность:
    \[
    P(A)=\frac{m}{n}.
    \]
\end{enumerate}

\section*{3. Очередность важна, дубли разрешены}
\sectionline

Пусть два человека независимо выбирают одну из цифр $1$, $2$, $3$.
Найдём вероятность того, что они выберут разные цифры.

Здесь:
\begin{itemize}
    \item строка --- выбор первого человека;
    \item столбец --- выбор второго человека;
    \item дубли разрешены: $11$, $22$, $33$ возможны;
    \item очередность важна: $12$ и $21$ --- разные исходы.
\end{itemize}

\begin{table}[ht]
\centering
\caption{Все исходы двух независимых выборов}
\begin{tabular}{@{}c c c c@{}}
\toprule
& \textbf{1} & \textbf{2} & \textbf{3} \\
\midrule
\textbf{1} & $11$ & \goodcell{$12$} & \goodcell{$13$} \\
\textbf{2} & \goodcell{$21$} & $22$ & \goodcell{$23$} \\
\textbf{3} & \goodcell{$31$} & \goodcell{$32$} & $33$ \\
\bottomrule
\end{tabular}
\end{table}

Всего допустимых клеток:
\[
n=9.
\]

Благоприятных клеток:
\[
m=6.
\]

Следовательно,
\[
P=\frac69=\frac23.
\]

\key{Ключевая мысль.}
Диагональ таблицы содержит случаи одинакового выбора:
\[
11,\qquad22,\qquad33.
\]

Все остальные клетки соответствуют выбору разных цифр.

\section*{4. Очередность не важна, дубли запрещены}
\sectionline

Пусть случайно выбирают две разные вершины квадрата $ABCD$ и соединяют
их отрезком. Найдём вероятность получить диагональ.

Здесь:
\begin{itemize}
    \item одну вершину нельзя выбрать дважды;
    \item очередность не важна: $AB$ и $BA$ задают один и тот же отрезок.
\end{itemize}

Поэтому диагональ таблицы исключаем, а из каждой пары зеркальных клеток
учитываем только одну.

Удобно оставить только верхнюю половину таблицы.

\begin{table}[ht]
\centering
\caption{Выбор двух различных вершин}
\begin{tabular}{@{}c c c c c@{}}
\toprule
& \textbf{$A$} & \textbf{$B$} & \textbf{$C$} & \textbf{$D$} \\
\midrule
\textbf{$A$}
& \badcell
& $AB$
& \goodcell{$AC$}
& $AD$ \\

\textbf{$B$}
& \badcell
& \badcell
& $BC$
& \goodcell{$BD$} \\

\textbf{$C$}
& \badcell
& \badcell
& \badcell
& $CD$ \\

\textbf{$D$}
& \badcell
& \badcell
& \badcell
& \badcell \\
\bottomrule
\end{tabular}
\end{table}

Допустимых исходов:
\[
n=6.
\]

Благоприятные исходы:
\[
AC,\qquad BD.
\]

Значит,
\[
m=2,
\qquad
P=\frac26=\frac13.
\]

\key{Типичная ошибка:}
одновременно учитывать $AB$ и $BA$.

Это один и тот же отрезок, поэтому такой подсчёт удваивает количество
исходов.

\section*{5. Случайная очередность трёх человек}
\sectionline

Пусть Андрей ($A$), Борис ($B$) и Виктор ($V$) случайно становятся
в очередь.

Найдём вероятность того, что Борис окажется позже Андрея.

Таблицу удобно строить так:
\begin{itemize}
    \item строка --- кто стоит первым;
    \item столбец --- кто стоит вторым;
    \item третий человек после выбора первых двух определяется автоматически;
    \item один человек не может занимать сразу два места, поэтому клетки
    главной диагонали запрещены.
\end{itemize}

\begin{table}[ht]
\centering
\caption{Все допустимые варианты очереди}
\begin{tabular}{@{}c c c c@{}}
\toprule
& \textbf{$A$ второй}
& \textbf{$B$ второй}
& \textbf{$V$ второй} \\
\midrule

\textbf{$A$ первый}
& \badcell
& \goodcell{$ABV$}
& \goodcell{$AVB$} \\

\textbf{$B$ первый}
& $BAV$
& \badcell
& $BVA$ \\

\textbf{$V$ первый}
& \goodcell{$VAB$}
& $VBA$
& \badcell \\

\bottomrule
\end{tabular}
\end{table}

Всего допустимых вариантов:
\[
n=6.
\]

Борис стоит позже Андрея в трёх случаях:
\[
ABV,\qquad AVB,\qquad VAB.
\]

Следовательно,
\[
m=3,
\qquad
P=\frac36=\frac12.
\]

\section*{6. Несколько испытаний с двумя исходами}
\sectionline

Пусть проводится три испытания.

В каждом испытании имеются два равновероятных исхода:
\[
+\qquad\text{и}\qquad -.
\]

Для трёх испытаний удобно:
\begin{itemize}
    \item результаты первых двух испытаний поместить в строки;
    \item результат третьего испытания поместить в столбцы.
\end{itemize}

\begin{table}[ht]
\centering
\caption{Три последовательных испытания}
\begin{tabular}{@{}c c c@{}}
\toprule
\textbf{Первые два}
& \textbf{Третий $+$}
& \textbf{Третий $-$} \\
\midrule
$++$ & $+++$ & \goodcell{$++-$} \\
$+-$ & \goodcell{$+-+$} & $+--$ \\
$-+$ & \goodcell{$-++$} & $-+-$ \\
$--$ & $--+$ & $---$ \\
\bottomrule
\end{tabular}
\end{table}

Пусть требуется найти вероятность ровно двух ``$+$''.

Благоприятны исходы:
\[
++-,\qquad +-+,\qquad -++.
\]

Всего клеток:
\[
n=8.
\]

Благоприятных:
\[
m=3.
\]

Поэтому
\[
P=\frac38=0{,}375.
\]

\key{Важно.}
Последовательности
\[
++-\qquad\text{и}\qquad +-+
\]
являются разными исходами.

Результат каждого испытания относится к своему номеру.

\section*{7. Выбор двух человек из группы}
\sectionline

Пусть из пяти туристов
\[
A,\ B,\ C,\ D,\ E
\]
случайно выбирают двоих.

Найдём вероятность того, что турист $A$ попадёт в выбранную пару.

Здесь:
\begin{itemize}
    \item одного человека нельзя выбрать дважды;
    \item очередность внутри пары не важна;
    \item поэтому учитываем только одну половину таблицы.
\end{itemize}

\begin{table}[ht]
\centering
\caption{Все допустимые пары туристов}
\begin{tabular}{@{}c c c c c c@{}}
\toprule
& $A$ & $B$ & $C$ & $D$ & $E$ \\
\midrule

$A$
& \badcell
& \goodcell{$AB$}
& \goodcell{$AC$}
& \goodcell{$AD$}
& \goodcell{$AE$} \\

$B$
& \badcell
& \badcell
& $BC$
& $BD$
& $BE$ \\

$C$
& \badcell
& \badcell
& \badcell
& $CD$
& $CE$ \\

$D$
& \badcell
& \badcell
& \badcell
& \badcell
& $DE$ \\

$E$
& \badcell
& \badcell
& \badcell
& \badcell
& \badcell \\

\bottomrule
\end{tabular}
\end{table}

Всего допустимых пар:
\[
n=10.
\]

С туристом $A$:
\[
AB,\qquad AC,\qquad AD,\qquad AE.
\]

Поэтому
\[
m=4,
\qquad
P=\frac4{10}=\frac25.
\]

\section*{8. Как понять, какую таблицу строить}
\sectionline

\begin{table}[ht]
\centering
\caption{Выбор структуры таблицы}
\begin{tabularx}{\linewidth}{@{}L{0.26\linewidth}L{0.29\linewidth}Y@{}}
\toprule
\textbf{Ситуация}
& \textbf{Как строить таблицу}
& \textbf{Что особенно проверить} \\
\midrule

Два независимых выбора
&
Первый выбор --- строки, второй --- столбцы
&
Разрешены ли одинаковые результаты; различаются ли $AB$ и $BA$
\\

Выбор двух разных объектов
&
Один объект --- строки, второй --- столбцы
&
Исключить диагональ; при неважной очередности оставить одну половину таблицы
\\

Очередь из трёх человек
&
Первое место --- строки, второе --- столбцы
&
На диагонали один человек занимал бы два места, поэтому она запрещена
\\

Три испытания с двумя исходами
&
Первые два результата --- строки, третий --- столбцы
&
Каждая последовательность должна встретиться ровно один раз
\\

Четыре испытания с двумя исходами
&
Первые два результата --- строки, последние два --- столбцы
&
Получится компактная таблица $4\times4$
\\

\bottomrule
\end{tabularx}
\end{table}

\section*{9. Главная схема решения}
\sectionline

Для большинства задач занятия достаточно следующего маршрута:

\[
\boxed{
\begin{array}{c}
\text{Допускаются ли дубли?}
\\[3pt]
\downarrow
\\[3pt]
\text{Важна ли очередность?}
\\[3pt]
\downarrow
\\[3pt]
\text{Что обозначают строки и столбцы?}
\\[3pt]
\downarrow
\\[3pt]
\text{Заполнить все допустимые клетки}
\\[3pt]
\downarrow
\\[3pt]
\text{Отметить благоприятные клетки}
\\[3pt]
\downarrow
\\[3pt]
P=\dfrac{m}{n}
\end{array}}
\]

\section*{10. Быстрая проверка таблицы}
\sectionline

Перед вычислением вероятности проверьте таблицу по шести вопросам.

\begin{table}[ht]
\centering
\caption{Самопроверка}
\begin{tabularx}{\linewidth}{@{}L{0.32\linewidth}Y@{}}
\toprule
\textbf{Вопрос}
& \textbf{Что проверить} \\
\midrule

Что означает одна клетка?
&
В каждой клетке записан один полностью определённый исход.
\\

Допустимы ли дубли?
&
Если повтор запрещён, соответствующие клетки исключены.
\\

Важна ли очередность?
&
Определено, считать ли $AB$ и $BA$ разными исходами.
\\

Нет ли двойного учёта?
&
При неважной очередности зеркальные клетки не считаются дважды.
\\

Все ли допустимые клетки заполнены?
&
Ни один возможный исход не пропущен.
\\

Все ли исходы равновозможны?
&
Только тогда можно непосредственно применять
\[
P=\frac{m}{n}.
\]
\\

\bottomrule
\end{tabularx}
\end{table}

\section*{11. Типичные ошибки}
\sectionline

\begin{itemize}
    \item
    \key{Посчитать зеркальные клетки дважды}, когда очередность не важна.

    \item
    \key{Удалить диагональ}, хотя дубли по условию разрешены.

    \item
    \key{Оставить диагональ}, хотя один объект нельзя выбрать повторно.

    \item
    \key{Считать $AB$ и $BA$ одинаковыми}, когда первый и второй выбор
    выполняют разные роли.

    \item
    \key{Пропустить исход} при бессистемном перечислении вариантов.

    \item
    \key{Смешать два способа подсчёта}: в знаменателе учитывать очередность,
    а в числителе её игнорировать.

    \item
    \key{Начать искать благоприятные случаи до заполнения всей таблицы.}

    Сначала строится пространство всех допустимых исходов, затем внутри него
    выделяется нужное событие.

    \item
    \key{Считать недопустимую клетку обычным исходом.}

    Перечёркнутая клетка вообще не входит в число $n$.
\end{itemize}

\clearpage

\section*{Тренировочный блок}
\sectionline

Во всех задачах используйте \key{табличный метод}.

Для каждой задачи:
\begin{enumerate}
    \item определите смысл строк и столбцов;
    \item решите, допускаются ли дубли;
    \item решите, важна ли очередность;
    \item заполните таблицу;
    \item отметьте благоприятные клетки;
    \item найдите $P=\dfrac{m}{n}$.
\end{enumerate}

\textbf{Средний уровень}

\begin{enumerate}[label=\textbf{\arabic*.}]
    \item
    Два ученика независимо выбирают по одной цифре из множества
    \[
    \{1,2,3\}.
    \]
    Найдите вероятность того, что они выберут разные цифры.

    \item
    У квадрата $ABCD$ случайно выбирают две различные вершины
    и соединяют их отрезком.
    Найдите вероятность того, что получится диагональ.

    \item
    Андрей, Борис и Виктор случайно становятся в очередь.
    Найдите вероятность того, что Виктор окажется раньше Андрея.
\end{enumerate}

\vspace{0.5em}

\textbf{Выше среднего}

\begin{enumerate}[label=\textbf{\arabic*.},resume]
    \item
    Два человека независимо выбирают по одному числу из множества
    \[
    \{1,2,3,4\}.
    \]
    Найдите вероятность того, что сумма выбранных чисел равна $5$.

    \item
    Из пяти туристов $A$, $B$, $C$, $D$, $E$ случайно выбирают двоих.
    Найдите вероятность того, что турист $A$ попадёт в выбранную пару.

    \item
    Монету бросают три раза.
    Каждый из двух результатов одного броска считается равновероятным.
    Найдите вероятность того, что орёл выпадет ровно два раза.

    \item
    Три друга $A$, $B$, $C$ случайно занимают три кресла подряд.
    Найдите вероятность того, что $A$ и $B$ будут сидеть рядом.

    \item
    У правильного пятиугольника $ABCDE$ случайно выбирают две различные
    вершины и соединяют их отрезком.
    Найдите вероятность того, что получится сторона пятиугольника.

    \item
    Последовательно выбирают две различные цифры из множества
    \[
    \{1,2,3,4\}.
    \]
    Первая выбранная цифра записывается в разряд десятков,
    вторая --- в разряд единиц.
    Найдите вероятность того, что первая цифра меньше второй.
\end{enumerate}

\vspace{0.5em}

\textbf{Сложный уровень}

\begin{enumerate}[label=\textbf{\arabic*.},resume]
    \item
    Монету бросают четыре раза.

    Для таблицы используйте строки, соответствующие результатам первых двух
    бросков, и столбцы, соответствующие результатам последних двух бросков.

    Найдите вероятность того, что орёл выпадет ровно два раза и при этом
    результаты первого и четвёртого бросков различны.
\end{enumerate}

\clearpage

\section*{Ответы}
\sectionline

\begin{table}[ht]
\centering
\caption{Ответы к тренировочному блоку}
\begin{tabularx}{0.78\linewidth}{@{}>{\bfseries}p{0.14\linewidth}X@{}}
\toprule
№ & Ответ \\
\midrule
1  & $\dfrac23$ \\
2  & $\dfrac13$ \\
3  & $\dfrac12$ \\
4  & $\dfrac14$ \\
5  & $\dfrac25$ \\
6  & $\dfrac38$ \\
7  & $\dfrac23$ \\
8  & $\dfrac12$ \\
9  & $\dfrac12$ \\
10 & $\dfrac14$ \\
\bottomrule
\end{tabularx}
\end{table}

\vfill

\begin{center}
\textcolor{softgray}{\small
Перед каждой задачей задайте три вопроса:
``допускаются ли дубли?'',
``важна ли очередность?'',
``что означает одна клетка таблицы?''
}
\end{center}

\end{document}